\documentclass[11pt,a4paper]{article}

\usepackage[utf8]{inputenc}
\usepackage[T1]{fontenc}
\usepackage[french]{babel}
\usepackage{lmodern}
\usepackage[margin=2cm]{geometry}
\usepackage{xcolor}
\usepackage{tikz}
\usepackage{booktabs}
\usepackage{tabularx}
\usepackage{array}
\usepackage{siunitx}
\usepackage{amsmath}
\usepackage{amssymb}
\usepackage{pifont}
\usepackage{enumitem}
\usepackage{listings}
\usepackage[colorlinks=true, linkcolor=blue!50!black, urlcolor=blue!60!black]{hyperref}

\usetikzlibrary{arrows.meta, positioning, shapes.geometric, calc, fit, backgrounds}

\definecolor{chw}{HTML}{B71C1C}
\definecolor{cdsp}{HTML}{1565C0}
\definecolor{ca53}{HTML}{2E7D32}
\definecolor{cusb}{HTML}{6A1B9A}
\definecolor{cnpu}{HTML}{E65100}
\definecolor{cfx}{HTML}{00838F}
\definecolor{cbuf}{HTML}{FFB300}
\definecolor{ccomp}{HTML}{455A64}
\definecolor{codebg}{HTML}{F5F5F5}

\tikzset{
  hw/.style={draw, fill=chw!15, rounded corners=2pt, font=\sffamily\footnotesize, align=center, minimum height=7mm, inner sep=3pt},
  dsp/.style={draw, fill=cdsp!15, rounded corners=2pt, font=\sffamily\footnotesize, align=center, minimum height=7mm, inner sep=3pt},
  a53/.style={draw, fill=ca53!15, rounded corners=2pt, font=\sffamily\footnotesize, align=center, minimum height=7mm, inner sep=3pt},
  pcm/.style={draw, fill=yellow!30, font=\sffamily\scriptsize, minimum height=5mm, rounded corners=1pt, align=center},
  fx/.style={draw, fill=cfx!15, rounded corners=2pt, font=\sffamily\footnotesize, align=center, minimum height=7mm, inner sep=3pt},
  usb/.style={draw, fill=cusb!15, rounded corners=2pt, font=\sffamily\footnotesize, align=center, minimum height=7mm, inner sep=3pt},
  arr/.style={-{Stealth[length=2mm]}, thick},
  arrmulti/.style={-{Stealth[length=2mm]}, line width=1.2pt}
}

\lstset{
  basicstyle=\ttfamily\footnotesize,
  backgroundcolor=\color{codebg},
  frame=single,
  rulecolor=\color{black!30},
  xleftmargin=0.5em,
  xrightmargin=0.5em,
  breaklines=true,
  columns=flexible,
  keepspaces=true,
  tabsize=2,
  inputencoding=utf8,
  extendedchars=true,
  literate={é}{{\'e}}1 {è}{{\`e}}1 {à}{{\`a}}1 {ê}{{\^e}}1 {â}{{\^a}}1 {ô}{{\^o}}1 {ù}{{\`u}}1 {û}{{\^u}}1 {î}{{\^\i}}1 {ï}{{\"\i}}1 {ç}{{\c c}}1 {É}{{\'E}}1 {È}{{\`E}}1 {À}{{\`A}}1 {→}{{$\to$}}1 {↓}{{$\downarrow$}}1
}

\title{\textbf{Plateforme Mastering Live Debix}\\[6pt]
  \Large Spécification Phase 1a --- \textbf{V3}\\[4pt]
  \large Console numérique 8 canaux end-to-end}
\author{Document de référence technique}
\date{2026-04-22 --- rev V3}

\begin{document}
\maketitle
\tableofcontents
\clearpage

% =======================================================================
\section{Contexte et reset}

\subsection{Historique V1 → V2.7 (résumé court)}

Les versions V1 à V2.7 (historique git complet dans \texttt{PHASE\_1A\_SPEC.pdf} commit 3da1cc21..8eac663c) ont tenté une architecture avec un flux stéréo séparé pour la master chain (\texttt{Source\_Play} 2ch) croisant un flux 8ch (\texttt{SAI\_Playback}). Le cross-pipeline mux demux a échoué au trigger IPC3 (\texttt{-110 timeout}).

\textbf{Leçon} : en SOF IPC3 sur i.MX8MP, un path playback asymétrique avec multiple cross-pipeline liens n'est pas fiable. Le pattern upstream qui fonctionne est tout-8ch linéaire.

\subsection{Nouvelle directive utilisateur}

Architecture \textbf{tout 8 canaux, end-to-end}. Le DSP traite 8 canaux de la même façon. Exports ALSA :
\begin{itemize}[leftmargin=*, itemsep=2pt]
  \item \texttt{SAI\_Capture} 8ch : les mics \textbf{post-effets IN}
  \item \texttt{SAI\_Playback} 8ch : entrée depuis DAW (ASIO return), injectée avant les effets OUT, après le mixer
  \item \texttt{HP\_Monitor} 8ch : la sortie finale \textbf{post-effets OUT} (pour NPU analyze)
  \item + USB gadget iOS \textbf{2IN/2OUT} séparé, entre dans le mixer
\end{itemize}

Flux interne DSP : mics → effets IN → routing 8×8 → MIXER (+ SAI\_Playback + iOS 2IN upmixé) → effets OUT → SAI TX + HP\_Monitor + iOS 2OUT downmixé.

\section{Architecture cible V3}

\begin{center}
\begin{tikzpicture}[node distance=4mm and 8mm, scale=0.85, transform shape]

\node[hw] (mics) {8 mics\\4× TAC5212 ADC};
\node[hw, below=14mm of mics] (hp) {HP/monitors\\4× TAC5212 DAC};

\node[dsp, right=16mm of mics, minimum width=22mm] (effin) {Effets IN\\8 canaux};
\node[pcm, right=10mm of effin] (cap) {SAI\_Capture\\ALSA 8ch cap};
\node[dsp, below=5mm of effin] (rout) {Routing 8×8};

\node[dsp, below=5mm of rout, minimum width=26mm] (mixer) {MIXER 8ch};
\node[pcm, right=10mm of mixer] (play) {SAI\_Playback\\ALSA 8ch play};
\node[usb, below=4mm of play] (ios) {iOS 2IN\\(USB gadget)};

\node[dsp, below=6mm of mixer, minimum width=22mm] (effout) {Effets OUT\\8 canaux};
\node[pcm, right=10mm of effout] (hpmon) {HP\_Monitor\\ALSA 8ch cap};
\node[usb, below=4mm of hpmon] (iosout) {iOS 2OUT\\(USB gadget)};

\draw[arrmulti] (mics.east) -- (effin.west);
\draw[arrmulti] (effin.east) -- (cap.west);
\draw[arrmulti] (effin.south) -- (rout.north);
\draw[arrmulti] (rout.south) -- (mixer.north);
\draw[arrmulti] (play.west) -- (mixer.east);
\draw[arrmulti] (ios.west) -| node[above, font=\tiny, pos=0.3]{upmix 2→8} (mixer.south east);
\draw[arrmulti] (mixer.south) -- (effout.north);
\draw[arrmulti] (effout.east) -- (hpmon.west);
\draw[arrmulti] (effout.south east) |- node[above, font=\tiny, pos=0.8]{downmix 8→2} (iosout.west);
\draw[arrmulti] (effout.west) -| (hp.north);

\end{tikzpicture}
\end{center}

\subsection{Les 4 PCMs ALSA (+ gadget USB)}

\begin{center}
\begin{tabularx}{\textwidth}{lllX}
\toprule
\textbf{PCM} & \textbf{Direction} & \textbf{Ch} & \textbf{Rôle} \\
\midrule
\texttt{SAI\_Capture} & capture & 8 & Mics post-effets IN (pour enregistrement DAW) \\
\texttt{SAI\_Playback} & playback & 8 & Retour DAW ASIO, entre dans le mixer \\
\texttt{HP\_Monitor} & capture & 8 & Sortie finale post-effets OUT (pour NPU) \\
USB gadget iOS & 2IN/2OUT & 2+2 & Séparé, connecte au mixer + reçoit downmix \\
\bottomrule
\end{tabularx}
\end{center}

\subsection{Le DSP applique toujours 8 canaux identiquement}

Un seul flux 8ch TDM de bout en bout. Les effets (compresseur multibande, volume, limiter) s'appliquent soit en \textbf{linked} (1 instance pour 8ch, même params) soit en \textbf{per-channel} (8 instances parallèles, params indépendants).

\section{Phasage réaliste}

\begin{center}
\begin{tabularx}{\textwidth}{llX}
\toprule
\textbf{Phase} & \textbf{Livrables} & \textbf{Objectif} \\
\midrule
\textbf{1a.1 MVP} & 2 pipelines 8ch basiques (capture effets IN + playback effets OUT), linked mode & Prouver que la base fonctionne \textbf{sans timeout trigger}. C'est ce qui a échoué en V2.x. \\
1a.2 & Matrice routing 8×8 + MIXER 8ch, injection SAI\_Playback & Pipeline playback complet avec retour DAW + routage mic \\
1a.3 & Tap HP\_Monitor (investiguer méthode : probe, DAI loop, ou autre) & Capture post-effets OUT pour NPU \\
1a.4 & USB gadget iOS 2IN/2OUT + upmix/downmix dans mixer & iOS connectable \\
2 & USB ASIO gadget 8IN/8OUT & DAW sur host externe \\
3 & Passage linked → per-channel (8 instances d'effets parallèles) & Paramètres d'effets indépendants par canal \\
\bottomrule
\end{tabularx}
\end{center}

\section{Phase 1a.1 MVP --- ce qu'on fait en premier}

\subsection{Architecture simplifiée}

\begin{center}
\begin{tikzpicture}[node distance=4mm and 8mm]

\node[hw] (mics) {mics\\(SAI RX 8ch)};
\node[dsp, right=10mm of mics, minimum width=40mm, fill=cdsp!20] (effin) {mbdrc 8ch → pga 8ch → drc 8ch};
\node[pcm, right=10mm of effin] (cap) {SAI\_Capture\\ALSA cap 8ch};

\node[pcm, below=22mm of mics] (play) {SAI\_Playback\\ALSA play 8ch};
\node[dsp, right=10mm of play, minimum width=40mm, fill=cdsp!20] (effout) {mbdrc 8ch → pga 8ch → drc 8ch};
\node[hw, right=10mm of effout] (hp) {SAI TX 8ch\\→ HP};

\draw[arrmulti] (mics) -- (effin);
\draw[arrmulti] (effin) -- (cap);
\draw[arrmulti] (play) -- (effout);
\draw[arrmulti] (effout) -- (hp);

\node[below=1mm of effin, font=\itshape\small]{Effets IN (linked, 1 instance × 8ch)};
\node[below=1mm of effout, font=\itshape\small]{Effets OUT (linked, 1 instance × 8ch)};

\end{tikzpicture}
\end{center}

\subsection{2 pipelines, aucun cross-pipeline}

\paragraph{PIPE 1 (capture avec effets IN)}
Extension directe de \texttt{pipe-drc-enabled-capture.m4} upstream : on enchaîne \texttt{mbdrc} + \texttt{pga} + \texttt{drc} en série, sur 8 canaux linked.

\paragraph{PIPE 6 (playback avec effets OUT)}
Nouvelle pipeline calquée sur \texttt{pipe-volume-playback.m4} mais avec \texttt{mbdrc} + \texttt{pga} + \texttt{drc} au lieu du seul \texttt{pga}. 8 canaux linked.

\paragraph{DAI\_ADD pour SAI7 RX+TX}
Identique à la baseline \texttt{drc8ms.m4}. Un DAI\_ADD capture + un DAI\_ADD playback.

\paragraph{PCMs exposés}
\texttt{PCM\_CAPTURE\_ADD(SAI\_Capture, 0, PIPELINE\_PCM\_1)} + \texttt{PCM\_PLAYBACK\_ADD(SAI\_Playback, 1, PIPELINE\_PCM\_6)}. 2 PCMs seulement en MVP (pas de \texttt{HP\_Monitor}, pas de \texttt{Master\_Tap}).

\subsection{Squelette m4 V3-A MVP}

\paragraph{\texttt{sof-imx8mp-tac5212-masterV3.m4} (top-level)}

\begin{lstlisting}[language={}, basicstyle=\ttfamily\scriptsize]
#
# Topology V3 Phase 1a.1 MVP : console 8ch end-to-end, linked effects.
# 2 pipelines 8ch sans cross-pipeline.
#
include(`utils.m4')
include(`dai.m4')
include(`pipeline.m4')
include(`sai.m4')
include(`pcm.m4')
include(`buffer.m4')
include(`common/tlv.m4')
include(`sof/tokens.m4')
include(`platform/imx/imx8.m4')

# PIPE 1 : SAI RX 8ch -> mbdrc + pga + drc -> SAI_Capture
PIPELINE_PCM_ADD(sof/sof-imx8mp-tac5212-pipe-effects-capture.m4,
    1, 0, 8, s32le,
    2000, 0, 0,
    48000, 48000, 48000)

# PIPE 6 : SAI_Playback -> mbdrc + pga + drc -> SAI TX
PIPELINE_PCM_ADD(sof/sof-imx8mp-tac5212-pipe-effects-playback.m4,
    6, 1, 8, s32le,
    2000, 0, 0,
    48000, 48000, 48000)

DAI_ADD(sof/pipe-dai-capture.m4,
    1, SAI, 7, tac5212-hifi,
    PIPELINE_SINK_1, 2, s32le,
    2000, 0, 0, SCHEDULE_TIME_DOMAIN_DMA)

DAI_ADD(sof/pipe-dai-playback.m4,
    6, SAI, 7, tac5212-hifi,
    PIPELINE_SOURCE_6, 2, s32le,
    2000, 0, 0, SCHEDULE_TIME_DOMAIN_DMA)

PCM_CAPTURE_ADD(SAI_Capture, 0, PIPELINE_PCM_1)
PCM_PLAYBACK_ADD(SAI_Playback, 1, PIPELINE_PCM_6)

DAI_CONFIG(SAI, 7, 0, tac5212-hifi,
    SAI_CONFIG(DSP_A, SAI_CLOCK(mclk, 12288000, codec_mclk_in),
        SAI_CLOCK(bclk, 12288000, codec_consumer),
        SAI_CLOCK(fsync, 48000, codec_consumer),
        SAI_TDM(8, 32, 255, 255),
        SAI_CONFIG_DATA(SAI, 7, 0)))
\end{lstlisting}

\paragraph{\texttt{pipe-effects-capture.m4} (structure)}
\begin{lstlisting}[language={}, basicstyle=\ttfamily\scriptsize]
# Capture : DAI SAI RX -> B0 -> mbdrc -> B1 -> pga -> B2 -> drc -> B3 -> HOST
# Pattern : pipe-drc-enabled-capture.m4 upstream étendu avec pga + drc.
include(`utils.m4','buffer.m4','pcm.m4','pipeline.m4',
        `bytecontrol.m4','mixercontrol.m4',
        `multiband_drc.m4','pga.m4','drc.m4','common/tlv.m4')

W_PCM_CAPTURE(PCM_ID, SAI Capture, 0, 2, SCHEDULE_CORE)

# Controls & blobs pour les 3 effets (canonique)
C_CONTROLMIXER(In Volume, PIPELINE_ID, ...)
define(DEF_PGA_TOKENS, ...) W_VENDORTUPLES(...) W_DATA(...)

include(`multiband_drc_coef_default.m4')
include(`drc_coef_default.m4')

W_MULTIBAND_DRC(0, PIPELINE_FORMAT, 2, 2, SCHEDULE_CORE, LIST(..))
W_PGA(0, PIPELINE_FORMAT, 2, 2, DEF_PGA_CONF, SCHEDULE_CORE, LIST(..))
W_DRC(0, PIPELINE_FORMAT, 2, 2, SCHEDULE_CORE, LIST(..))

W_BUFFER(0..3, 8ch config)

P_GRAPH(pipe-effects-capture, PIPELINE_ID,
    LIST(`dapm(N_BUFFER(0), N_DAI_IN)',
         `dapm(N_MULTIBAND_DRC(0), N_BUFFER(0))',
         `dapm(N_BUFFER(1), N_MULTIBAND_DRC(0))',
         `dapm(N_PGA(0), N_BUFFER(1))',
         `dapm(N_BUFFER(2), N_PGA(0))',
         `dapm(N_DRC(0), N_BUFFER(2))',
         `dapm(N_BUFFER(3), N_DRC(0))',
         `dapm(N_PCMC(PCM_ID), N_BUFFER(3))'))

indir(`define', concat(`PIPELINE_SINK_', PIPELINE_ID), N_BUFFER(0))
indir(`define', concat(`PIPELINE_PCM_', PIPELINE_ID), SAI Capture PCM_ID)
PCM_CAPABILITIES(SAI Capture PCM_ID, ...)
\end{lstlisting}

\paragraph{\texttt{pipe-effects-playback.m4} (structure)}
Même squelette que capture mais inversé : \texttt{W\_PCM\_PLAYBACK} en entrée, exports \texttt{PIPELINE\_SOURCE} en sortie pour le DAI TX.

\subsection{Tests Phase 1a.1 MVP}

\begin{enumerate}[leftmargin=*, itemsep=2pt]
  \item \textbf{T0 Load} : kernel charge la tplg sans \texttt{-22} ni \texttt{-110}. Carte \texttt{softac5212tdm} + 2 PCMs visibles (\texttt{SAI\_Capture}, \texttt{SAI\_Playback}).
  \item \textbf{T1 Playback OK} : \texttt{aplay -D plughw:softac5212tdm,1 tone.wav} → son audible sur HP, pas de xrun.
  \item \textbf{T2 Capture OK} : \texttt{arecord -D plughw:softac5212tdm,0 -d 3 mics.wav} → 8 canaux enregistrés, signal sur les mics branchés.
  \item \textbf{T3 Full-duplex} : aplay + arecord simultanés, pas de xrun ni désync.
  \item \textbf{T4 Effets audibles} : changer \texttt{amixer cset name='Out Volume'} pendant playback → volume change audible et sans glitch.
  \item \textbf{T5 Preset mbdrc} : modifier le blob \texttt{Master MBDRC Config} → compression audible changer.
\end{enumerate}

Si T0-T5 passent, on passe à Phase 1a.2 (ajout mixer + routing).

\section{Phase 1a.2 --- routing 8×8 + MIXER}

Une fois le MVP stable, on injecte le routage :

\begin{itemize}[leftmargin=*, itemsep=1pt]
  \item Composant \texttt{mux} en mode routing 8→8 après les effets IN (patch bay mics)
  \item Composant \texttt{mixer} 8ch combinant : mics routés + SAI\_Playback + iOS 2IN upmixé
  \item Impact topology : plus de canaux, plus de composants, mais toujours en-ligne (pas de cross-pipeline)
\end{itemize}

Design détaillé à écrire en V3.2 après validation MVP.

\section{Phase 1a.3 --- tap HP\_Monitor (point dur)}

Le tap post-effets OUT pour NPU est le vrai défi architectural.

Options à investiguer :
\begin{enumerate}[leftmargin=*, itemsep=1pt]
  \item \textbf{Mux demux en fin de pipeline playback} : celle qui a échoué en V2.7 (trigger timeout). À ré-investiguer dans le contexte V3 simplifié.
  \item \textbf{SOF Probes} : a été invalidé côté driver kernel i.MX8MP (investigation 919e09d7).
  \item \textbf{Pipeline capture synchronisée} : pipeline capture supplémentaire qui lit le buffer de sortie des effets OUT via DAPM. Pattern jamais testé pour playback.
  \item \textbf{Loopback hardware} : câble court d'une sortie TAC vers un ADC TAC libre. Pattern studio pro (SSL/Neve "insert meter").
\end{enumerate}

À décider empiriquement après Phase 1a.1.

\section{Questions à l'investigation critic}

\begin{enumerate}[leftmargin=*, itemsep=2pt]
  \item \textbf{Linked vs per-channel} : les composants \texttt{multiband\_drc}, \texttt{drc}, \texttt{pga} dans notre fork SOF (branche \texttt{feature/sdma-ap2ap-phase1}) traitent-ils les 8 canaux en mode linked (1 envelope/gain global) ou per-channel (état par canal) ? Réponse factuelle par lecture de \texttt{src/audio/multiband\_drc/multiband\_drc\_generic.c} et \texttt{drc\_generic.c}.
  \item \textbf{pipe-drc-enabled-capture.m4 étendu} : peut-on sérier \texttt{multiband\_drc → pga → drc} dans une pipeline capture en IPC3 sans blocage ? Quelles contraintes (buffer sizes, core, periods) à respecter ?
  \item \textbf{pipe-volume-playback.m4 étendu} : idem pour la pipeline playback. Ajouter \texttt{multiband\_drc} et \texttt{drc} avant le \texttt{pga} existant — est-ce trivial ou des patterns upstream particuliers requis ?
  \item \textbf{8ch end-to-end vs contraintes plateforme} : \texttt{PLATFORM\_MAX\_CHANNELS = 8} sur i.MX8MP confirmé. Les composants \texttt{multiband\_drc}, \texttt{drc}, \texttt{pga} supportent-ils tous 8 canaux en entrée/sortie ?
  \item \textbf{Kconfig requis} : au-delà de \texttt{CONFIG\_COMP\_MULTIBAND\_DRC=y} et \texttt{CONFIG\_COMP\_VOLUME=y} déjà activés en V2.x, y a-t-il d'autres configs à ajouter ?
  \item \textbf{Tap post-effets OUT en IPC3 playback} : quelle est la méthode canonique (hors smart-amp qui est capture) pour extraire le signal post-effets d'une pipeline playback vers un PCM capture ALSA ? Existe-t-il un précédent IPC3 upstream ?
  \item \textbf{Pourquoi V2.7 a échoué au trigger} : avec la nouvelle vision claire, est-ce que l'hypothèse "cross-pipeline playback non supporté" est correcte, ou l'échec venait d'autre chose (blob malformé, chain widgets, etc.) ?
\end{enumerate}

\section{Livrables Phase 1a.1 MVP}

\begin{itemize}[leftmargin=*, itemsep=1pt]
  \item \texttt{sof-imx8mp-tac5212-masterV3.m4} : top-level topology
  \item \texttt{sof-imx8mp-tac5212-pipe-effects-capture.m4} : pipeline capture 8ch avec effets
  \item \texttt{sof-imx8mp-tac5212-pipe-effects-playback.m4} : pipeline playback 8ch avec effets
  \item Board config \texttt{imx8mp\_evk\_mimx8ml8\_adsp.conf} : déjà à jour (V2.3)
  \item Script de test \texttt{test-phase1a1-mvp.sh}
\end{itemize}

\section{Protocole}

Aucune modification SOF sans \texttt{critic\_analyze} préalable. Un changement = un commit = validation utilisateur. Gate compile-test \texttt{m4 | alsatplg} OBLIGATOIRE avant deploy.

\end{document}
